%=====================================================================
%  v6 Supplementary Information -> LaTeX
%=====================================================================
\documentclass[11pt,a4paper]{article}
\usepackage[a4paper,top=2.2cm,bottom=2.2cm,left=2.0cm,right=2.0cm]{geometry}
\usepackage{fontspec}
\setmainfont{Times New Roman}
\setsansfont{Arial}
\setmonofont{Consolas}
\usepackage{amsmath,amssymb}
\usepackage{booktabs,longtable,array,tabularx}
\usepackage{graphicx}
\newcommand{\figfit}[2][\textwidth]{%
  \includegraphics[width=#1,height=0.82\textheight,keepaspectratio]{#2}}
\usepackage{pdflscape}
\usepackage{threeparttable}
\usepackage[super,square,sort&compress]{natbib}
\usepackage[hidelinks]{hyperref}
\usepackage{caption}
\captionsetup{font=small,labelfont=bf}
\setcounter{figure}{0}
\renewcommand{\thefigure}{S\arabic{figure}}
\setcounter{table}{0}
\renewcommand{\thetable}{S\arabic{table}}
\usepackage{microtype}
\usepackage{parskip}

\graphicspath{{../figures_v5/}{./}}

\title{\bfseries Supplementary Information\\[6pt]
\large Disentangling Between-Target and Within-Target Signal in
PPI $\Delta\Delta G$ Benchmarking}
\author{}
\date{}

\begin{document}
\maketitle

\section*{Supplementary Note 1 | Derivation of the covariance decomposition}
Write the true value and prediction for mutation $i$ in target $c$ as
$y_{ci} = \mu_c + \varepsilon_{ci}$ and
$\hat{y}_{ci} = \hat{\mu}_c + \hat{\varepsilon}_{ci}$. Expanding
$\operatorname{Cov}(y,\hat{y})$ gives a between-target term and a within-target term,
\begin{equation}
  \operatorname{Cov}(y,\hat{y})
  = \frac{1}{N}\sum_{c} n_c \left(\mu_c - \bar{\mu}\right)
    \left(\hat{\mu}_c - \bar{\hat{\mu}}\right)
  + \frac{1}{N}\sum_{c}\sum_{i} \varepsilon_{ci}\,\hat{\varepsilon}_{ci};
\end{equation}
all cross terms vanish identically because within-target deviations have zero mean in every
target. Dividing by $\sigma_y \sigma_{\hat{y}}$ gives $C_{\mathrm{B}}$ and
$C_{\mathrm{W}}$ and establishes
\begin{equation}
  r_{\mathrm{pooled}} = C_{\mathrm{B}} + C_{\mathrm{W}}
\end{equation}
exactly. The identity residual is reported as a quality-control quantity (observed maximum
$3.331 \times 10^{-16}$). Because $C_{\mathrm{B}}$ and $C_{\mathrm{W}}$ share the same
denominator they are covariance-signal contributions on a common scale, not independent
correlations.

\section*{Supplementary Note 2 | Row universes and their relationship}
The controlled-split universe (7,523 rows, 230 complexes,
112 structural families) is the full set of forward-paired true
$\Delta\Delta G$ observations. It contains 4,812 distinct
(complex, mutation) keys; 2,011 keys carry more than one
experimental record and 1,405 of those disagree in
$\Delta\Delta G$.\par
The annotation namespace (4,473 rows, 214 targets,
108 structural families) is the feature-complete subset of that universe:
every one of its 4,473 rows is present in the controlled-split universe, and
3,050 rows of the latter are absent from it
(16 targets and
4 structural families are not
represented). The two row counts therefore differ because they describe different
filtering stages of the same underlying data, not different benchmarks.\par
The target-count distributions are:
\begin{center}\begin{tabular}{@{}lrrrrr@{}}\toprule
Universe & $n \geq 1$ & $n \geq 2$ & $n \geq 3$ & $n \geq 5$ & $n \geq 10$ \\
\midrule
Controlled-split (7,523) & 230 & 203 & 193 & 172 & 118 \\
Annotation (4,473) & 214 & 180 & 164 & 130 & 85 \\
\bottomrule\end{tabular}\end{center}

\section*{Supplementary Note 3 | Matched universes and row-set fingerprints}
Each prediction output received a SHA-256 fingerprint of its sorted row identifiers.
Cross-model comparison was performed only between outputs sharing a fingerprint; raw metrics
were never compared across fingerprints.

\begin{landscape}
\begin{table}[h]
  \centering
  \scriptsize
  \begin{threeparttable}
  \caption{Matched-evaluation strata and row-set fingerprints.}
  \label{tab:s1}
  \begin{tabular}{@{}lllrrrrrrr@{}}
    \toprule
    Universe & Artifact & Row-set fingerprint & Rows ($N$) & Targets & Families
      & Targets ($n \geq 3$) & Targets ($n \geq 5$) & Targets ($n \geq 10$) & Models \\
    \midrule
U\_V6\_2\_split4\_8d003e00 & V6\_2\_split4 & 8d003e00057d932a & 7523 & 230 & 112 & 193 & 172 & 118 & 36 \\
U\_V6\_2\_1\_donly\_8d003e00 & V6\_2\_1\_donly & 8d003e00057d932a & 7523 & 230 & 112 & 193 & 172 & 118 & 6 \\
U\_V7.1\_z\_34b45937 & V7.1\_z & 34b45937b42b7114 & 4473 & 214 & 108 & 164 & 130 & 85 & 10 \\
U\_V7.0\_ablation\_34b45937 & V7.0\_ablation & 34b45937b42b7114 & 4473 & 214 & 108 & 164 & 130 & 85 & 35 \\
U\_V7.0\_losses\_34b45937 & V7.0\_losses & 34b45937b42b7114 & 4473 & 214 & 108 & 164 & 130 & 85 & 16 \\
U\_V7.1\_lsr\_34b45937 & V7.1\_lsr & 34b45937b42b7114 & 4473 & 214 & 108 & 164 & 130 & 85 & 6 \\
U\_V7.2R1\_full\_2f5a5c2e & V7.2R1\_full & 2f5a5c2e4832298a & 4233 & 174 & 91 & 142 & 116 & 75 & 12 \\
U\_V7.2R2\_trackB\_e81801ff & V7.2R2\_trackB & e81801ff29e01b7b & 4219 & 174 & 91 & 142 & 115 & 75 & 12 \\
U\_V7.2R3\_ppi\_17487417 & V7.2R3\_ppi & 174874178847c9bc & 3452 & 132 & 70 & 105 & 87 & 55 & 15 \\
U\_V7.2R3\_noise\_17487417 & V7.2R3\_noise & 174874178847c9bc & 3452 & 132 & 70 & 105 & 87 & 55 & 12 \\
U\_V7.3\_physics\_f8c33c7e & V7.3\_physics & f8c33c7e987d72c8 & 3351 & 166 & 86 & 129 & 102 & 65 & 30 \\
U\_V7.2\_nested\_c08fe062 & V7.2\_nested & c08fe062e64e4cba & 1807 & 13 & 9 & 13 & 12 & 11 & 12 \\
    \bottomrule
  \end{tabular}
  \end{threeparttable}
\end{table}
\end{landscape}

\section*{Supplementary Note 4 | Derivation of fixed-$K$ shortlist discriminability}
To separate candidate-set geometry from the realised magnitudes of $\Delta\Delta G$, we
evaluate fixed-$K$ selection against a \emph{canonical rank-linear utility reference}.
Rank the $n$ candidates of a target from most to least favourable and assign the $j$-th
ranked candidate the reference utility
\begin{equation}
  u_{(j)} = n - j + 1, \qquad j = 1,\dots,n,
  \label{eq:rankutil}
\end{equation}
so that utilities are equally spaced and strictly decreasing in rank. This reference is a
deliberate idealisation: it fixes the \emph{shape} of the utility profile and asks how much
of the total utility a size-$K$ shortlist can capture as a function of $n$ and $K$ alone.
It is not a model of any particular $\Delta\Delta G$ distribution, and the resulting
quantity is a geometric surrogate for shortlist discriminability rather than a property of
real affinity data. Let $E = \sum_{i \in S_K} u_i / \sum_{i=1}^{n} u_i$ denote the share of
total utility captured by a selected set $S_K$ of $K$ candidates. Two references define the
discriminable interval:
\begin{itemize}
  \item $E_{\max}$ --- the share attained by the optimal (oracle) selection. Under
        Eq.~\eqref{eq:rankutil} the top $K$ candidates have utilities
        $n, n-1, \dots, n-K+1$, and $\sum_{j=1}^{n} u_{(j)} = n(n+1)/2$, so
        \begin{equation}
          E_{\max} = \frac{\sum_{j=1}^{K}(n-j+1)}{n(n+1)/2}
                    = \frac{K\,(2n-K+1)}{n\,(n+1)} .
          \label{eq:emax}
        \end{equation}
        The closed form depends on the reference Eq.~\eqref{eq:rankutil}; for an arbitrary
        utility vector $E_{\max} = \sum_{j\le K} u_{(j)} / \sum_{j\le n} u_{(j)}$ and is
        not determined by $n$ and $K$ alone.
  \item $E_{\mathrm{rand}}$ --- the share expected under uniform random selection of $K$ of the
        $n$ candidates. By symmetry of the uniform distribution over subsets, each candidate
        is selected with probability $K/n$ irrespective of its utility, so
        \begin{equation}
          E_{\mathrm{rand}} = \frac{K}{n} .
          \label{eq:erand}
        \end{equation}
        This reference is exact under any utility profile, including the rank-linear one.
\end{itemize}
The discriminable interval is their difference, which simplifies in closed form:
\begin{equation}
  D(n,K) \;=\; E_{\max} - E_{\mathrm{rand}}
  \;=\; \frac{K(2n-K+1)}{n(n+1)} - \frac{K}{n}
  \;=\; \frac{K\,(n-K)}{n\,(n+1)} .
  \label{eq:dnk}
\end{equation}
For $n \leq K$ every candidate is selected, hence $E_{\max} = E_{\mathrm{rand}} = 1$ and
$D = 0$ identically: such targets carry no information for fixed-$K$ shortlist metrics.
Treating $n$ as continuous and differentiating Eq.~\eqref{eq:dnk},
\begin{equation}
  \frac{\mathrm{d}D}{\mathrm{d}n}
  = K\,\frac{(n+1) - (n-K)(2n+1)}{n^{2}(n+1)^{2}} = 0
  \;\Longrightarrow\; n - K = \frac{n(n+1)}{2n+1},
\end{equation}
whose positive root is
\begin{equation}
  n^{*} = K + \sqrt{K^{2} + K} .
  \label{eq:nstar2}
\end{equation}
For $K = 10$ this gives $n^{*} \approx 20.5$, with $D(20) = 0.2381$ and
$D(50) = 0.1569$; in the present benchmark $53$ of $130$ targets with $n \geq 5$ ($40.8\%$)
fall in the degenerate regime $n \leq K$.

\emph{Applicability and limits.} Everything above is conditional on the canonical
rank-linear reference Eq.~\eqref{eq:rankutil}. $D(n,K)$ is therefore a geometric
surrogate for shortlist discriminability, not a statement about arbitrary utility
distributions: with any other utility profile $E_{\max}$ changes and so does $D$.
It is not a model-performance metric, it does not measure the $\Delta\Delta G$ utility range
of any real dataset, and it does not imply that SKEMPI $\Delta\Delta G$ values are
intrinsically similar. Establishing that would require an analysis using realised utility
magnitudes, which we did not perform. What is reference-free is the degenerate case: for
$n \leq K$ a size-$K$ shortlist selects every candidate, so $E_{\max} = E_{\mathrm{rand}} = 1$
and $D = 0$ under \emph{any} utility profile. That statement, and the target counts it
implies, does not depend on Eq.~\eqref{eq:rankutil}.

\section*{Supplementary Note 5 | Weighted covariance decomposition}
The five target-weighting schemes act on the \emph{target level}: predictions are unchanged and
only the weight assigned to each target varies. Let $n_c$ be the number of mutations in target
$c$. Because the schemes are named after target-aggregation rules, we take the target weight
$A_c$ as the quantity they describe,
\begin{equation}
  A_c \in \{\, n_c,\; 1,\; \sqrt{n_c},\; \log(1+n_c),\; \min(n_c, 30) \,\}
  \label{eq:Ac}
\end{equation}
for the per-mutation, per-target, $\sqrt{n}$, logarithmic and capped schemes respectively, and
split it evenly across the target's rows:
\begin{equation}
  w_{ci} = \frac{A_c}{n_c}, \qquad
  \sum_{i} w_{ci} = A_c, \qquad
  A = \sum_{c}\sum_{i} w_{ci} = \sum_{c} A_c .
  \label{eq:wci}
\end{equation}
This is the only reading under which the scheme names are accurate: with $A_c = 1$ each target
carries total weight $1$ regardless of its size (genuinely per-target), and with $A_c = n_c$
every row carries weight $1$ (genuinely per-mutation). Writing $w_{ci} = a_c$ instead --- the
literal form present in the original implementation --- would give each target total weight
$n_c a_c$, so that the ``per-target'' scheme would still let large targets dominate; we
therefore treat that literal reading as a sensitivity variant
(Eq.~\eqref{eq:sens} below).

Under the measure in Eq.~\eqref{eq:wci} the weighted means and variances are
\begin{equation}
  \bar{y}_w = \frac{1}{A}\sum_{c}\sum_{i} w_{ci}\, y_{ci}, \qquad
  \sigma^{2}_{y,w} = \frac{1}{A}\sum_{c}\sum_{i} w_{ci}\,(y_{ci} - \bar{y}_w)^{2},
  \label{eq:means}
\end{equation}
and likewise for $\hat{y}$. Because $w_{ci} = A_c/n_c$ is constant within a target, the
within-target mean is the ordinary target mean,
\begin{equation}
  B_{y,c} \;=\; \frac{\sum_i w_{ci}\, y_{ci}}{\sum_i w_{ci}} \;=\; \bar{y}_c .
  \label{eq:byc}
\end{equation}
With $\varepsilon_{ci} = y_{ci} - B_{y,c}$ and
$\hat{\varepsilon}_{ci} = \hat{y}_{ci} - B_{\hat{y},c}$, the components are
\begin{align}
  C_{\mathrm{B},w} &= \frac{\sum_{c} A_c\,
      \bigl(B_{y,c} - \bar{y}_w\bigr)\bigl(B_{\hat{y},c} - \bar{\hat{y}}_w\bigr)}
      {A\,\sigma_{y,w}\,\sigma_{\hat{y},w}}, \label{eq:cbw}\\[4pt]
  C_{\mathrm{W},w} &= \frac{\sum_{c}\sum_{i} w_{ci}\,
      \varepsilon_{ci}\,\hat{\varepsilon}_{ci}}
      {A\,\sigma_{y,w}\,\sigma_{\hat{y},w}}. \label{eq:cww}
\end{align}
Substituting $y_{ci} = B_{y,c} + \varepsilon_{ci}$ and
$\hat{y}_{ci} = B_{\hat{y},c} + \hat{\varepsilon}_{ci}$ into
$A\,\sigma_{y,w}\sigma_{\hat{y},w}\,r_w = \operatorname{Cov}_w(y,\hat{y})$ gives
\begin{equation}
  \sum_{c}\sum_{i} w_{ci}\,(y_{ci} - \bar{y}_w)(\hat{y}_{ci} - \bar{\hat{y}}_w)
  = \underbrace{\sum_{c} A_c
      \bigl(B_{y,c} - \bar{y}_w\bigr)\bigl(B_{\hat{y},c} - \bar{\hat{y}}_w\bigr)}_{
      \text{Eq.~\eqref{eq:cbw}}}
  + \underbrace{\sum_{c}\sum_{i} w_{ci}\,\varepsilon_{ci}\hat{\varepsilon}_{ci}}_{
      \text{Eq.~\eqref{eq:cww}}},
  \label{eq:cross}
\end{equation}
because for each target $\sum_i w_{ci}\varepsilon_{ci} = 0$ and
$\sum_i w_{ci}\hat{\varepsilon}_{ci} = 0$, so both cross terms vanish. Hence
\begin{equation}
  r_w \;=\; C_{\mathrm{B},w} + C_{\mathrm{W},w}
  \label{eq:identity-w}
\end{equation}
holds exactly for every scheme and every evaluation instance. Verified numerically on all 109
configurations crossed with all five schemes (545 instances): maximum absolute residual
$5.6 \times 10^{-17}$.

\emph{Sensitivity to the weighting reading.} If instead the row weight is taken literally as
$w_{ci} = a_c$ with $a_c$ given by Eq.~\eqref{eq:Ac}, the target total weight becomes
$n_c a_c$ and the same algebra applies with
\begin{equation}
  w^{\mathrm{sens}}_{ci} = a_c, \qquad \sum_i w^{\mathrm{sens}}_{ci} = n_c\,a_c .
  \label{eq:sens}
\end{equation}
The identity still holds exactly, because numerator and denominator share one measure. The two
readings differ only in \emph{what a scheme equalises}, and therefore in the size of the
reported components:

\begin{center}
\begin{tabular}{@{}lrr@{}}
\toprule
Quantity & Eq.~\eqref{eq:wci} (primary) & Eq.~\eqref{eq:sens} (sensitivity) \\
\midrule
Median cross-scheme range of $C_{\mathrm{B}}$ & 0.3052 & 0.0930 \\
Median cross-scheme range of $C_{\mathrm{W}}$ & 0.1240 & 0.0624 \\
$C_{\mathrm{B}}$ sign reversals (any magnitude) & 40/109 & 44/109 \\
$C_{\mathrm{B}}$ sign reversals with $|C_{\mathrm{B}}|>0.005$ & 34/109 & 32/109 \\
$C_{\mathrm{W}}$ sign reversals (any magnitude) & 5/109 & 0/109 \\
\bottomrule
\end{tabular}
\end{center}

Both readings agree on the qualitative conclusion: the between-target component is markedly
less stable across weighting conventions than the within-target component (a larger
cross-scheme range and many more sign reversals), so a reported decomposition must state its
weighting scheme. They disagree on whether $C_{\mathrm{W}}$ ever reverses sign
(5/109 under the primary reading versus 0/109 under the sensitivity reading),
and on the absolute magnitude of the ranges; the main text quotes the primary reading and
reports the sensitivity here.

\section*{Supplementary Note 6 | Duplicate-key sensitivity}
The controlled-split universe contains 7{,}523 mutation records but only 4{,}812 distinct
(complex, mutation) keys: 2{,}011 keys carry more than one experimental record, and 1{,}405 of
those disagree in $\Delta\Delta G$. This raises a legitimate concern: could part of the
random-split performance, and hence part of the apparent collapse under family-out-of-domain
evaluation, be an artefact of the same key appearing more than once in the evaluation set?

Because the stored prediction artifacts do not retain per-row fold assignments, the split
cannot be regenerated here; we therefore quantify the sensitivity \emph{within} the stored
evaluation set by recomputing every quantity under three key treatments:
\begin{itemize}
  \item \textbf{ALL} --- all 7{,}523 records (the reported analysis);
  \item \textbf{FIRST} --- one record per key, keeping the first occurrence (4{,}812 records);
  \item \textbf{MEAN} --- one record per key, with $y$ and $\hat{y}$ both averaged within the
        key (4{,}812 records).
\end{itemize}
The decomposition used here is the unweighted identity of Supplementary Note 1, and was
verified to reproduce the frozen per-configuration values to 3.3e-16 across the 36
controlled-split configurations.

\begin{center}
\begin{tabular}{@{}llrrrr@{}}
\toprule
Regime & Key treatment & Median $r$ & Median $C_{\mathrm{B}}$
  & Median $C_{\mathrm{W}}$ & Median macro $\rho$ \\
\midrule
Random & ALL & +0.6844 & $+0.4275$ & $+0.2591$ & 0.0707 \\
Random & FIRST & $+0.6378$ & $+0.3920$ & $+0.2465$ & 0.0886 \\
Random & MEAN & $+0.6338$ & $+0.3705$ & $+0.2647$ & 0.0977 \\
\midrule
Family-OOD & ALL & +0.1508 & $-0.0072$ & $+0.1580$ & 0.0377 \\
Family-OOD & FIRST & $+0.1614$ & $-0.0122$ & $+0.1759$ & 0.0354 \\
Family-OOD & MEAN & $+0.2181$ & $+0.0197$ & $+0.1956$ & 0.0347 \\
\bottomrule
\end{tabular}
\end{center}

The paired Random $\to$ Family-OOD differences are:
\begin{center}
\begin{tabular}{@{}lrrrr@{}}
\toprule
Key treatment & $\Delta r$ & $\Delta C_{\mathrm{B}}$ & $\Delta C_{\mathrm{W}}$
  & Between-target share \\
\midrule
ALL   & -0.5332 & -0.4324 & -0.1011 & 81.0\,\% \\
FIRST & -0.4733 & -0.4027 & -0.0706 & 85.1\,\% \\
MEAN  & -0.4266 & -0.3476 & -0.0671 & 83.8\,\% \\
\bottomrule
\end{tabular}
\end{center}

Three observations follow. First, collapsing duplicate keys \emph{reduces} the random-split
correlation (+0.6844 to $+0.6378$ under FIRST) while slightly \emph{increasing} the
family-out-of-domain correlation (+0.1508 to $+0.1614$); the collapse is therefore, if
anything, understated rather than inflated by repeated records. Second, $\Delta C_{\mathrm{B}}$
remains large and negative under every treatment, and its share of the summed absolute
component change lies between 81.0\,\% and 85.1\,\%: the qualitative conclusion that the
between-target component carries most of the decline does not depend on how repeated keys are
handled. Third, the magnitude of the decline is sensitive to the treatment
(-0.5332 versus -0.4266), so the absolute size of the collapse should be quoted with the key
treatment stated.

\emph{Limits.} This is an evaluation-set-internal sensitivity analysis. It cannot rule out
duplicate-key leakage across the train/test boundary, because that requires regenerating the
splits under key-level grouping and retraining, which the stored artifacts do not permit and
which we did not perform. The 1{,}405 keys with disagreeing $\Delta\Delta G$ values also
indicate that a (complex, mutation) key does not always identify a unique experimental
condition; determining whether those conditions should enter the key definition is a
data-curation decision that we did not make.

\section*{Supplementary Note 7 | External coverage and common intersection}
The five external methods are annotation columns of the same published SKEMPI table, mapped to
the annotation namespace through the normalised key of
\texttt{scripts/v8\_1/05\_external\_clean.py}. Each method has its own pattern of missing
entries, so the uniquely mapped subsets differ slightly in size: Flex-ddG maps 1{,}755 rows,
DDGPred 1{,}756, and FoldX, ESM-IF and proteinMPNN 1{,}776, out of 1{,}781 uniquely mappable
rows. Because all five columns originate from the same annotation universe and share common
row identifiers, we can remove the coverage imbalance by restricting every method to
their common row intersection.

\begin{center}
\begin{tabular}{@{}lrrrrrrr@{}}
\toprule
Method & Rows & Targets ($n \geq 5$) & Pooled Pearson & Macro Spearman
  & $C_{\mathrm{B}}$ & $C_{\mathrm{W}}$ & RMSE \\
\midrule
\multicolumn{8}{@{}l@{}}{\emph{(A) variable coverage --- each method's own uniquely mapped subset, raw columns}} \\
DDGPred & 1,756 & 65 & +0.2406 & +0.1287 & +0.0715 & +0.1690 & 2.1262 \\
ddG\_FoldX & 1,776 & 66 & +0.2089 & +0.1088 & +0.0657 & +0.1433 & 2.3691 \\
Flex-ddG & 1,755 & 65 & +0.2028 & +0.1092 & +0.0478 & +0.1550 & 2.1636 \\
proteinMPNN & 1,776 & 66 & +0.1509 & +0.0320 & +0.0699 & +0.0810 & --- \\
ESM-IF & 1,776 & 66 & -0.2250 & +0.0213 & -0.1149 & -0.1101 & --- \\
\midrule
\multicolumn{8}{@{}l@{}}{\emph{(A$'$) same subsets, exploratory sign-oriented values}} \\
DDGPred & 1,756 & 65 & +0.2406 & +0.1287 & +0.0715 & +0.1690 & 2.1262 \\
ddG\_FoldX & 1,776 & 66 & +0.2089 & +0.1088 & +0.0657 & +0.1433 & 2.3691 \\
Flex-ddG & 1,755 & 65 & +0.2028 & +0.1092 & +0.0478 & +0.1550 & 2.1636 \\
proteinMPNN & 1,776 & 66 & +0.1509 & +0.0320 & +0.0699 & +0.0810 & --- \\
ESM-IF & 1,776 & 66 & +0.2250 & -0.0213 & +0.1149 & +0.1101 & --- \\
\midrule
\multicolumn{8}{@{}l@{}}{\emph{(B) common intersection --- raw columns}} \\
DDGPred & 1,752 & 65 & +0.2410 & +0.1287 & +0.0714 & +0.1696 & 2.1265 \\
ddG\_FoldX & 1,752 & 65 & +0.2121 & +0.0967 & +0.0681 & +0.1440 & 2.3623 \\
Flex-ddG & 1,752 & 65 & +0.2036 & +0.1092 & +0.0476 & +0.1560 & 2.1620 \\
proteinMPNN & 1,752 & 65 & +0.1497 & +0.0430 & +0.0690 & +0.0807 & --- \\
ESM-IF & 1,752 & 65 & -0.2225 & +0.0099 & -0.1113 & -0.1112 & --- \\
\midrule
\multicolumn{8}{@{}l@{}}{\emph{(B$'$) common intersection, exploratory sign-oriented values}} \\
DDGPred & 1,752 & 65 & +0.2410 & +0.1287 & +0.0714 & +0.1696 & 2.1265 \\
ddG\_FoldX & 1,752 & 65 & +0.2121 & +0.0967 & +0.0681 & +0.1440 & 2.3623 \\
Flex-ddG & 1,752 & 65 & +0.2036 & +0.1092 & +0.0476 & +0.1560 & 2.1620 \\
proteinMPNN & 1,752 & 65 & +0.1497 & +0.0430 & +0.0690 & +0.0807 & --- \\
ESM-IF & 1,752 & 65 & +0.2225 & -0.0099 & +0.1113 & +0.1112 & --- \\
\bottomrule
\end{tabular}
\end{center}

RMSE is shown only for Flex-ddG, FoldX and DDGPred, whose columns are on a $\Delta\Delta G$-like
scale; ESM-IF and proteinMPNN return dimensionless mutation-effect scores for which an error in
kcal/mol is not defined (Supplementary Note~8).

The common intersection retains 1,752 of the 1{,}781 uniquely mappable rows (98.4\,\%),
covering 163 targets of which 65 have $n \geq 5$. Comparing the treatments:

\begin{itemize}
  \item on the \emph{raw} columns the two rankings are strongly positively aligned
        (Spearman $\rho = +0.90$, exact permutation $p = 0.083$), as in the main text;
  \item after the exploratory sign-orientation the same ranks show no significant agreement
        ($\rho = +0.30$, $p = 0.683$), and ESM-IF moves from last to second on the
        pooled axis while remaining last on the within-target axis;
  \item in both treatments DDGPred ranks first on both axes.
\end{itemize}

The coverage imbalance is therefore not what drives the difference between the two conclusions:
restricting all five methods to their common intersection leaves the raw panel aligned and the
oriented panel unaligned, exactly as in the full uniquely mapped subsets. The panel remains small
(five methods) and underpowered under either treatment.

\section*{Supplementary Note 8 | Provenance, orientation and units of the external predictions}
The five external methods are not separate downloads: they are five annotation columns of the
SKEMPI-derived protein--protein interaction table distributed with CATH-ddG (main text
reference~3; \texttt{skempi\_v2.csv}), which we read directly without modification
(Table~S3). This matters for interpretation in three ways.

\emph{The published columns are preserved unmodified.} For every method we ingest the
published column exactly as distributed. We apply no affine rescaling and no recalibration,
and the raw values reproduce to four decimals when recomputed directly from the published
columns on the same uniquely mapped subset (e.g.\ ESM-IF $-0.2250$ and ProteinMPNN
$+0.1509$). The only operation applied anywhere is the single sign multiplier $o_m$
introduced in the main text Methods (\emph{Score orientation and its limits}). It is used
solely for the exploratory
orientation sensitivity and never modifies the raw columns reported here; because $o_m$ is
selected from the evaluation data, that sensitivity is explicitly label-informed, whereas
the raw analysis requires no such choice and is the primary result. So the reported raw
correlations are properties
of the published scores and not of any post-processing.

\emph{Orientation.} The analysed response is $y = \Delta\Delta G = \Delta G_{\mathrm{mut}} -
\Delta G_{\mathrm{WT}}$ (Methods), so $y > 0$ denotes weaker binding. We use every column
exactly as published and apply no sign inversion to any method. Two consequences are worth
stating. First, the Flex-ddG, FoldX and DDGPred columns are already oriented so that larger
values mean weaker binding, matching $y$, and they correlate positively with it. Second, the
published ESM-IF column is oriented the other way --- larger values mean stronger binding ---
which is why it carries a negative pooled correlation ($-0.2250$). That negative sign is a
property of the published ESM-IF score on this benchmark, not an artefact of our processing;
we did not flip it, and we flag it here so that the sign is not misread.

\emph{Units, and why two RMSE values are not reported.} Flex-ddG, FoldX and DDGPred are
quantitative $\Delta\Delta G$-like predictors (signed binding-affinity changes) whose
columns span standard deviations of
1.31, 1.92 and 1.55 respectively, so an error in kcal/mol is meaningful for them. ESM-IF and
ProteinMPNN are inverse-folding sequence models: their published outputs are dimensionless
mutation-effect scores with standard deviations of 0.093 and 0.012, roughly two orders of
magnitude narrower than the experimental $\Delta\Delta G$ spread. An RMSE expressed in kcal/mol
is therefore not defined for those two columns, and we report ``---''. This follows the rule
already applied elsewhere in this work, where a dimensionless mutation-effect score is recorded
with \texttt{RMSE = NOT\_APPLICABLE} rather than converted into an apparent energy error.
Correlation-based quantities --- pooled Pearson, macro within-target Spearman,
$C_{\mathrm{B}}$ and $C_{\mathrm{W}}$ --- are invariant to positive affine rescaling and are reported for all five methods.

\emph{Sign orientation for cross-method comparison.} The published columns are not all in
the same direction. Using within-complex correlations against $y = \Delta\Delta G$ --- which
removes between-target differences and therefore isolates direction rather than predictive
skill --- Flex-ddG, FoldX, DDGPred and ProteinMPNN are aligned with $y$ ($+0.49$, $+0.49$,
$+0.51$ and $+0.30$), whereas ESM-IF is anti-aligned ($-0.24$): its published score is larger
for stronger binding. For the cross-method performance comparison we therefore apply a
sign-only orientation, $s^{\mathrm{aligned}} = o_m\, s^{\mathrm{raw}}$ with
$o_{\mathrm{ESM\text{-}IF}} = -1$ and $o_m = +1$ for the other four (Methods). This uses no
evaluation labels and no recalibration, and it changes no information content: it only flips
the sign of the correlation-like quantities. Table~S3 records each raw column with its raw
orientation and $o_m$, so the unmodified provenance and the aligned values are both
auditable. The effect on the panel is material: used raw the two axes appear strongly
aligned ($\rho = +0.90$, $p = 0.083$), whereas after orientation they show no significant
agreement ($\rho = +0.30$, $p = 0.683$). We report the oriented values throughout and state
the raw alternative explicitly, since with five methods neither reaches adequate power.

\emph{Scope of the comparison.} Because the two score-based methods are not on the same scale
as the three energy-based ones, the RMSE column in Table~3 should not be used to compare across
those groups; the pooled and within-target correlations are the quantities that are comparable
across all five methods, and they are the basis of the rank-agreement analysis.


\section*{Machine-readable supplementary data}
The full model-family and configuration registry --- including per-configuration row-set
fingerprints, split regimes, training-provenance flags, ablation and representative markers,
and the representative-selection reason --- is supplied as machine-readable Supplementary
Data (\texttt{MODEL\_FAMILY\_REGISTRY.csv}, \texttt{MATCHED\_UNIVERSE\_REGISTRY.csv},
\texttt{FROZEN\_INPUT\_MANIFEST.csv}). The weighting audit configurations are enumerated in
\texttt{WEIGHTING\_FINAL.csv}, and the duplicate-key sensitivity in
\texttt{DUPLICATE\_KEY\_SENSITIVITY.csv}.

\section*{Supplementary Figures}
\begin{landscape}\begin{table}[h]\centering\footnotesize
\begin{threeparttable}
\caption{Model pools used in the rank-agreement sensitivity analysis.}\label{tab:s4pools}
\begin{tabular}{@{}llrrrr@{}}\toprule
Row-set hash & Level & Models & Method families & Green models & $\rho$(pooled, macro) \\
\midrule
f8c33c7e & A & 30 & 10 & 30 & +0.8242 \\
8d003e00 & A & 42 & 12 & 12 & +0.7596 \\
8d003e00 & C & 12 & 12 & 3 & +0.6923 \\
c08fe062 & A & 12 & 4 & 12 & +0.5315 \\
2f5a5c2e & A & 12 & 4 & 12 & +0.4336 \\
e81801ff & A & 12 & 4 & 12 & +0.3427 \\
17487417 & C & 9 & 9 & 9 & +0.0769 \\
17487417 & A & 27 & 9 & 27 & -0.0174 \\
34b45937 & C & 34 & 34 & 27 & -0.0716 \\
34b45937 & A & 67 & 34 & 60 & -0.2685 \\
\bottomrule\end{tabular}
\begin{tablenotes}\footnotesize
\item \emph{Level A} ($\ast$LA in Figure~S3) denotes a provenance-clean stratum:
every listed model is free of known supervised training overlap with the evaluation
rows. \emph{Level C} ($\ast$LC) denotes a permissive-split stratum, where the pool
includes configurations trained under random/complex/\texttt{seq40} partitions whose
evaluation rows may overlap their training folds; these are retained only for this
sensitivity analysis. `Green models' is the number of models in the pool flagged
provenance-clean. Pools within the same row-set hash are comparable because they
share identical evaluation rows; pools with different hashes are not.
\end{tablenotes}\end{threeparttable}\end{table}\end{landscape}
\begin{landscape}\begin{table}[h]\centering\scriptsize
\begin{threeparttable}
\caption{Provenance and score semantics of the external prediction columns.}\label{tab:s3external}
\begin{tabular}{@{}llllll@{}}\toprule
Method & Raw column & Units & Raw orientation & $o_m$ & RMSE meaningful? \\
\midrule
Flex-ddG & \texttt{Flex-ddG} & kcal/mol (ΔΔG-like scaling) & higher = larger $\Delta\Delta G$ (weaker binding) & $+1$ & yes \\
DDGPred & \texttt{DDGPred} & kcal/mol (ΔΔG-like scaling) & higher = larger $\Delta\Delta G$ (weaker binding) & $+1$ & yes \\
ddG\_FoldX & \texttt{ddG\_FoldX} & kcal/mol (ΔΔG-like scaling) & higher = larger $\Delta\Delta G$ (weaker binding) & $+1$ & yes \\
ESM-IF & \texttt{ESM-IF} & dimensionless (no kcal/mol scaling) & higher = stronger binding & $-1$ & \textbf{no} \\
proteinMPNN & \texttt{proteinMPNN} & dimensionless (no kcal/mol scaling) & higher = larger $\Delta\Delta G$ (weaker binding) & $+1$ & \textbf{no} \\
    \bottomrule\end{tabular}
\begin{tablenotes}\footnotesize
\item All five columns come from one published annotation table of SKEMPI-derived
protein--protein interaction mutations. Raw values are used exactly as published; we apply
no affine calibration and no recalibration against the evaluation labels. The $o_m$ column
records the sign multiplier used only for the exploratory orientation sensitivity.
We verified this by recomputing every reported correlation directly from the raw columns,
which reproduces the reported values to four decimals for all five methods. Raw standard
deviations differ by more than an order of magnitude across columns (1.31, 1.55 and 1.92
for the three $\Delta\Delta G$-like predictors versus 0.093 and 0.012 for the two
inverse-folding scores), which is why an RMSE in kcal/mol is defined only for the former.
\end{tablenotes}\end{threeparttable}\end{table}\end{landscape}

\begin{figure}[htbp]
  \centering
  \figfit{FigS1_matched_universe_strata_EN.pdf}
  \caption{Matched-evaluation strata grouped by SHA-256 row-set
  fingerprints. Each bar is one prediction output; outputs sharing a fingerprint are
  directly comparable with each other and with no others.}
\end{figure}

\begin{figure}[htbp]
  \centering
  \figfit{FigS2_capability_heatmap_provenance_clean_EN.pdf}
  \caption{Capability profiles of matched, provenance-clean ablation
  variants (seed 42). Rows are controlled ablation variants of one local-physics method
  family, not independent methods. Values are standardised within each metric for visual
  comparison only; raw values are printed in each cell.}
\end{figure}

\begin{figure}[htbp]
  \centering
  \figfit{FigS3_pool_sensitivity_rank_agreement_EN.pdf}
  \caption{Pool sensitivity of cross-metric rank agreement across
  matched strata. Each bar is one model pool within one row-set stratum, labelled
  \texttt{artifact (row-set hash / level)}. \emph{LA} denotes a provenance-clean stratum
  (no known supervised training overlap) and \emph{LC} a permissive-split stratum retained
  only for this sensitivity analysis; pool sizes, method-family counts and inclusion
  criteria are listed in Table~S2 below. Within one row-set hash the evaluation rows are
  identical and only the set of included model families changes, which moves the agreement
  between pooled Pearson and macro within-target Spearman from $-0.26$ to $+0.82$; this
  quantity is therefore reported as a range rather than a single value.}

\end{figure}

\end{document}
